\documentclass[UTF8,fontset=fandol,a4paper,zihao=-4]{ctexart}
\usepackage[a4paper,margin=2.5cm]{geometry}
\usepackage{hyperref,booktabs,array,enumitem}
\hypersetup{hidelinks,pdftitle={AI工具使用详情},pdfauthor={}}
\setlist{itemsep=0.3em,parsep=0pt}
\title{AI工具使用详情}
\author{}\date{}
\begin{document}
\maketitle
\noindent\textbf{材料状态：研究复核版。官方模拟器测试按要求搁置；没有参赛队实际完成逐项人工核验的记录。本文件不构成签名、合规承诺或可直接提交的证明。}

\section{工具名称及可核实的型号信息}
当前主对话模型标识为\texttt{mcm/gpt-6-astra}，通过Oh My Pi（OMP）编码工作流调用；历史阶段的模型信息应以实际会话记录为准，不推测供应方内部版本。并行子代理承担几何、覆盖、环境、策略候选、代码审查、数学审查、图表及支撑材料任务；未向本记录提供的子代理后端型号不作臆测。

Python、NumPy、SciPy、Matplotlib、PyMuPDF及XeLaTeX是实际计算、制图和编译工具，不被当作独立AI模型。Python统一在conda的\texttt{py314}环境运行。实际依赖版本由\path{results/environment_manifest.json}记录，复现命令见\path{reproduce.txt}。

\section{具体使用目的与参与环节}
\begin{enumerate}
\item\textbf{审题与模型建立：}AI读取B题及两份附件，抽取频道唯一性、示向误差、距离边界、方向半平面、光学条件与计时规则，提出四问之间的统一模型。
\item\textbf{核心数学推导：}AI参与半平面交、区域直径、可实现的等边三角形反例、最小包围圆判据、第二测点三圆交候选域，以及七点和31点不漏检覆盖证明。这些均属核心建模，不能表述为仅做语言润色。
\item\textbf{代码和算法：}AI编写几何计算、观测驱动规划器、局部丢信号恢复、光学覆盖、本地合成环境、正式HTTP客户端、实验与验证脚本。
\item\textbf{计算和结果：}AI调用本机程序实际运行试验、核对逐动作计时、整理数据及图表；没有由语言模型凭空填写数值。
\item\textbf{独立自动审查：}未参与对应原实现的审查代理依据题面和计算证据检查模型与代码。它仍是AI审查，不是参赛队员人工审核。
\item\textbf{论文与交付：}AI撰写正文、附录、图注、复现说明及本文件，并安排LaTeX编译和PDF逐页检查。最终检查是否完成，以随包证据记录为准，不以本声明代替。
\end{enumerate}

\section{主要提示方式与过程}
用户的主要要求依次为：尽可能详细、高质量作答2026年B题；所有Python使用conda的\texttt{py314}；按照目录中\texttt{prompt.md}要求完成从问题、模型、算法、结果、验证到论文的闭环；后来明确说明官方模拟器目前不能运行，要求搁置。

\texttt{prompt.md}中的关键约束包括：LLM负责思考、本机负责可程序确定的计算；先建立简单可运行方案；在统一接口下进行独立探索和审查；增强须有同口径对照；不得虚构实验、资料或日志；编译成功不能替代PDF视觉检查。上述要求被转为分阶段任务和实际代码接口。

对子代理采用明确的目标、文件所有权和验收条件，例如：证明未知方向下的不漏检；只通过公开动作响应决策而不读取隐藏源；所有数值由主流程统一执行验证；未运行的组件不得自称已经通过；不得把合成结果称为正式成绩。主要工作顺序为：原题与规则核实、几何和覆盖组件、合成环境及客户端、基线与共享策略、独立审查、定向失信号恢复开发、隔离种子的最终试验、图表生成、论文整合及交付检查。

\section{AI输出的采纳、修改与自动核验}
本项目采纳了保守定位集合、几何覆盖骨架及逐频道完成证书作为主线，没有采用没有证据的误差独立性、按场强反推距离或猜测总数停止等做法。重要自动纠正记录如下：
\begin{itemize}
\item “直径不超过40米即可一次清除”被明确否定，以最小包围圆半径不超过20米为正确判据。
\item 同方差信息矩阵不直接称作Fisher信息；仅在显式高斯辅助似然下使用该名称，主模型仍用有界误差。
\item 零误差示向度的背向射线漏洞通过显式正投影约束修复，保存了回归检查。
\item 过小格距和几乎无数值余量的格距被限制，默认31点构造增加正方向余量证明。
\item HTTP截断响应按同ID恢复，持续慢速响应按绝对时间期限失败关闭，不推进未确认状态。
\item 基线、批量扫描和共享观测分别对照；不把全部收益归给共享，也不宣称每个实例都更优。
\end{itemize}

核心计算证据包括：
\begin{itemize}
\item\path{results/final_experiments.json}：4430次最终合成运行。
\item\path{results/enhancement_experiments.json}：3080次独立种子的增强比较与压力检验。
\item 新增环收缩解析覆盖、31个格点不可删数值证据、静态2-opt和自适应测点负结果、滚动2-opt以及完整终止状态回归。
\item\path{results/ring_sweep.json}：8个环半径开发比较与500个新世界配对确认，共2200次运行；逐世界结果可复算，未将已看过的检验组冒充新确认组。
\item\path{results/geometry_verification.json}：独立包围圆、包含性及覆盖检查。
\item\path{results/intersection_verification.json}：有界、无界及空集实例。
\item 两份HTTP故障检查JSON、第二测点示例、恢复策略开发与最终消融。
\end{itemize}
结果文件明确区分数学证明、数值采样和本地合成环境。最终支撑包还记录逐页视觉检查及解压复现结果。

\section{人工修改和人工核验情况}
\textbf{尚无实际参赛队员逐项人工核验记录。}本次会话中，用户确认了选B题、Python环境、工作要求以及官方模拟器暂缓；这些决策不是对每条模型、代码和结果的人工审查。程序修改、结果核对及多代理审查均不能被记作人工完成。

如后续用于参赛，参赛队应实际阅读并理解核心证明与源代码，复算关键结果，检查模型是否符合题意，审阅全部AI参与内容并据实记录采纳、人工修改和核验情况。不能直接将本节改写为“已全部核验”，除非对应工作确实完成。

\section{官方测试与提交边界}
曾查阅官方说明及下载目录，但未成功运行官方模拟器，未产生官方演练成绩、六次正式案例编码或六份加密导出日志。用户明确要求搁置后，不再继续下载、启动或调用官方模拟器。本地合成JSON和HTTP夹具日志不能替代这些材料。

官方要求使用AI时在论文参考文献之前声明，并在支撑材料中提供名为“AI工具使用详情.pdf”的文件；未经必要人工审查的核心AI成果不能直接作为合规参赛成果提交。

规定来源：全国大学生数学建模竞赛组委会，《人工智能工具使用规定（2026年试行）》。
\begin{flushleft}\small
\url{https://www.mcm.edu.cn/html_cn/node/fef94648f2836ab6cc81586f4c38512b.html}
\end{flushleft}

本研究材料没有伪造参赛队号、个人签名、审核记录、测试日志、上传状态或获奖预测。含参赛队号的未来私有HTTP日志不得未经脱敏进入匿名支撑包；官方加密日志则须按要求保留原名及原始字节。
\end{document}
